\chapter{Fluids}

The central concepts for fluids are 
\begin{itemize}[noitemsep]
\item derivation of the governing equations of motion 
in the current configuration,
\item a stress response function 
that is decomposed into volumetric and
deviatoric parts,
\item a volumetric stress response function 
that is either incompressible
or compressible, and 
\item a deviatoric stress response that, 
through dependence on velocity gradients, 
is either free from shear (inviscid) or 
supports shear (viscid).
\end{itemize}

\section{Incompressible, Inviscid (Ideal) Fluid}
The incompressible, inviscid (ideal) has three
assumptions:
\be
\begin{tabular}{cl|l}
 &
 {\bf Assumption} 
 & 
 {\bf Expression}  \\ \hline \hline
 (1) 
 & 
 The reference mass density is uniform. 
 & 
 $\rho_0(\vX) = \bar\rho_0 > 0$ 
 \\ \hline
 (2) 
 &
 The material is incompressible.
 & 
 $\grad \cdot \vv(\vx, t) = 0$
 \\ \hline
 (3)
 &
 The Cauchy stress is spherical only.
 & 
 \multirow{ 2}{*}{$\vsigma(\vx, t) = - P(\vx, t) \; \vid$} \\
 
 &
 The stress has no deviatoric response.
 & 
\end{tabular}
\label{eq:incompressible_inviscid}
\ee

\begin{Hresult}
The current mass density $\rho$ is a positive constant, equal to the
reference mass density.  
Thus $\rho(\vx, t) = \rho_0(\vX) = \bar\rho_0 > 0$.
\end{Hresult}

\begin{Hproof}
Use conservation of mass in Eulerian form 
(\ref{eq:mass_conservation}),
\be
\frac{\partial \rho(\vx, t)}{\partial t} + \grad \cdot [
\rho(\vx, t) \; \vv(\vx, t) ] = 0.
\ee
Expand the divergence operator $\grad \cdot [ \; \cdot \; ]$ on the 
product of density and velocity,
\be
\underset{=\dot\rho \;\; (\ref{eq:total_time_derivative_scalar})}{\underbrace{
\frac{\partial \rho(\vx, t)}{\partial t} 
+ [ \grad \rho(\vx, t) ] \; \vv(\vx, t)  
}}
+  \rho(\vx, t) \; 
\underset{=0 \;\; (\ref{eq:incompressible_inviscid})}
{\underbrace{\grad \cdot \vv(\vx, t)}}
= 0.
\ee
Thus,
\be
\dot{\rho}(\vx, t) = 0,
\ee
and expressed in the reference configuration variable $\vX$,
\be
\frac{\partial \rho(\vvarphi(\vX, t), t)}{\partial t} = 0.
\ee
Since the time variation is zero at the initial time $t=0$,
\be
\rho(\vvarphi(\vX, 0), 0) 
= \rho(\vX, 0) = \rho_0(\vX) = \bar\rho_0,
\ee
\be
\therefore \; \rho(\vx, t) =  \bar\rho_0 > 0.
\label{eq:rho_x_t_rho0}
\ee
Thus the current mass density $\rho(\vx, t)$ 
is uniform and constant for all time.
\end{Hproof}

\begin{Hresult}
The {\bf Euler equations} for an {\bf ideal fluid} can then be written as
\be
\boxed{
\bar\rho_0 
\left[ 
\frac{\partial \vv}{\partial t} 
+ 
\left[\grad \vv\right]
\; \vv
\right] 
= -\grad P + \bar\rho_0 \; \vb
\;\; \Longleftrightarrow \;\;
\bar\rho_0 
\left[
 v_{i,t} + v_{i,j} \; v_j
\right]
= 
- P_{,i} 
+
\bar\rho_0 \;
b_i
}
\ee
with
\be
\boxed{
\grad \cdot \vv = 0
\;\; \Longleftrightarrow \;\;
v_{i,i} = 0
}
\ee
where functional dependence on the current configuration, 
$f = f(\vx, t)$, 
is assumed, but not explicitly written for the sake of brevity.
\end{Hresult}

\begin{Hproof}
The appearance of $\rho(\vx, t)$ in the current configuration balance
of linear momentum is substituted with $\bar\rho_0$, 
\be
\bar\rho_0 \;\dot{\vv}(\vx, t) 
= \grad \cdot \vsigma(\vx, t) + \bar\rho_0 \; \vb(\vx, t).
\ee
The $\grad \cdot \vsigma$ term 
reduces to $\grad \cdot (-P \vid) = - \grad P$.
The spatial body force per unit mass is $\vb(\vx, t)$.
The material time derivative is written explicitly as
\be
\dot{\vv}(\vx, t) 
= \frac{\partial \vv(\vx, t)}{\partial t} 
+ [\grad \vv(\vx, t)] \; \vv(\vx, t).
\ee
The Euler equations for an ideal fluid can then be written as
\be
\bar\rho_0 \left[ 
\frac{\partial \vv(\vx, t)}{\partial t} 
+ 
\left[\grad \vv(\vx, t)\right] \; \vv(\vx, t)
\right] 
= -\grad P(\vx, t) + \bar\rho_0 \; \vb(\vx, t), 
\ee
with the incompressibility constraint, 
\be
\grad \cdot \vv(\vx, t)= 0.
\ee
used to obtain the foregoing result.
\end{Hproof}

\begin{Hremark}
\begin{itemize}[noitemsep]
\item The ideal fluid does not support shear stresses.
\item There is no explicit relationship between pressure $P$ and 
density $\bar\rho_0$.
\item There is no explicit relationship between pressure $P$ and 
velocity $\vv$.
\end{itemize}
\end{Hremark}

\section{Compressible, Inviscid Fluid}
The compressible, inviscid fluid eliminates the constant density and 
incompressibility
assumptions of an ideal fluid.
To accommodate the compressibility behavior, the pressure becomes a function
of density, $P = P(\rho(\vx, t), t)$, with the added constraint that
changes in pressure with respect to changes in density must be positive.
Figure~\ref{fig:P_rho_diagram} 
illustrates this relationship
(compare to Figure~\ref{fig:PV_diagram}).
\be
\begin{tabular}{cl|l}
 &
 {\bf Assumption} 
 & 
 {\bf Expression}  \\ \hline \hline
 (1) 
 & 
 The Cauchy stress is spherical only, 
 & 
 \multirow{ 3}{*}{$\vsigma(\vx, t) = - P(\rho(\vx, t),t) \; \vid$} \\
 
 &
 and is a function of density. 
 & 
 \\
 
 &
 The stress has no deviatoric response.
 & 
 \\ \hline
 (2)
 & 
 The change in pressure with respect to  
 &
 \multirow{ 2}{*}{$\dst\frac{\partial P}{\partial \rho} > 0$}
 \\
 &
 change in density is positive.
 & 
\end{tabular}
\label{eq:compressible_inviscid}
\ee


\begin{figure}[htb]
  \begin{center}
  %\begin{overpic}[width=0.7\textwidth, grid, tics=10]{fig/P_rho_diagram}
  \begin{overpic}[width=0.7\textwidth,natwidth=165,natheight=97]{fig/P_rho_diagram.pdf}
    \put( 4.5, 56){$P$}
    \put(96,  7){$\rho$}
    \put(54,  2){$\rho_0$}
    \put( 0, 18.5){$P_0$}
    \put(65, 13){slope $=\dst\frac{\partial P}{\partial \rho} > 0$}
    \put( 8,  2){expansion}
    \put(66,  2){compression}
  \end{overpic}
  \end{center}
  \caption{Pressure as a function of density for a compressible, inviscid
  fluid.}
  \label{fig:P_rho_diagram}
\end{figure}

\begin{Hremark}
Refer to Figure~\ref{fig:P_rho_diagram}.
We discuss the $P(\rho)$ relationship
near a reference density $\rho_0$, at reference pressure $P_0$.
Near $\rho_0$, the slope is
positive, increases as $\rho$ increases, 
and decreases and $\rho$ decreases.
The slope remains positive for
all values of $\rho$, which is a positive, real value.

The character of $P(\rho)$ at the extreme values of $\rho$, zero and
infinity, are shown in general terms only.  The density of the fluid
can tend toward zero, but as the pressure tends toward zero (a vacuum), 
the fluid will disassociate and the continuity assumption underlying
continuum mechanics ceases to be valid.  Thus, the $P(\rho)$ curve is
plotted with a hole where the $P$ and $\rho$ axes intersect.

Conversely, as the density of the fluid tends toward infinity, so too 
does the pressure.  While no vertical density asymptote is shown
in Figure~\ref{fig:P_rho_diagram}, it is implied because density 
increasing without bounds seems nonphysical.  The sharp upward 
bend of $P(\rho)$ suggests this asymptotic behavior.

The general behavior that the slope of the $P(\rho)$ curve
always be positive follows from the notion that an increase in pressure
on a fluid should cause increase in density, and vice versa.  Thus, 
the slope of the curve should never be zero or negative.

Finally, unlike the first derivative $\partial P / \partial \rho > 0$, 
the second derivative 
$\partial^2 P / \partial \rho^2$ 
has no requirements of positivity, even though 
$P(\rho)$ is shown
in Figure~\ref{fig:P_rho_diagram} as a convex function.
\end{Hremark}

\begin{Hresult}
The {\bf elastic fluid equations} can then be written as
\be
\boxed{
\rho
\left[ 
\frac{\partial \vv}{\partial t} 
+ 
\left[\grad \vv\right]
\; \vv
\right] 
= -\dst\frac{\partial P}{\partial \rho} \grad \rho + \rho \; \vb
\;\; \Longleftrightarrow \;\;
\rho
\left[
 v_{i,t} + v_{i,j} \; v_j
\right]
= 
- \dst\frac{\partial P}{\partial \rho} \rho_{,i} 
+
\rho \;
b_i
}
\label{eq:elastic_fluid_momentum}
\ee
with
\be
\boxed{
\dst\frac{\partial \rho}{\partial t} 
+
\grad \cdot
( \rho \; \vv) 
= 
0
\;\; \Longleftrightarrow \;\;
\rho_{,t}
+ 
(\rho \; v_{i})_{,i} = 0
}
\label{eq:elastic_fluid_mass}
\ee
where functional dependence on the current configuration, 
$f = f(\vx, t)$, 
is assumed, but not explicitly written for the sake of brevity.
\end{Hresult}

\begin{Hremark}
\begin{itemize}[noitemsep]
\item The elastic fluid is compressible, and thus allows
volume and density changes.
\item The elastic fluid does not support shear stresses.
\item There is an explicit relationship between pressure $P$ and 
density $\rho$, and the first derivative of this relationship must
be greater than zero.  
\item There is no explicit relationship between pressure $P$ and 
velocity $\vv$.
\item The elastic fluid is the starting point for the discussion of
shock waves in condensed matter (see Chapter~\ref{chapter:shock}).
\end{itemize}
\end{Hremark}


\subsection{The Wave Equation}
Now we add additional assumptions to the assumptions of a 
compressible, inviscid fluid, listed in 
(\ref{eq:compressible_inviscid}).


\be
\begin{tabular}{cl|l}
 &
 {\bf Assumption} 
 & 
 {\bf Expression}  \\ \hline \hline
 (1) 
 & 
 Body forces are negligible compared 
 & 
 \multirow{ 2}{*}{$\vb(\vx, t) = \ve{0}$} 
 \\
 % 
 &
 to internal and inertial forces.
 & 
 \\ 
 \hline
 (2)
 & 
 Change in pressure with respect to change 
 &
 \multirow{ 2}{*}{$\dst\frac{\partial P}{\partial \rho} = c^2$}
 \\
 &
 in density is a positive (squared) constant.
 & 
 \\
 \hline
 (3)
 & 
 Uniform density as in (\ref{eq:rho_x_t_rho0})
 &
 $\rho(\vx,t) = \bar\rho_0$
 \\
 \hline
 (4)
 & 
 Quiescent initial velocity 
 &
 $\vv(\vx,0) = \ve{0}$
 \\
 \hline
 (5)
 & 
 The velocity field is a perturbation $\epsilon$
 &
 \multirow{ 2}{*}{$\vv(\vx,t) = \epsilon \; \vv(\vx,0)$}
 \\
 % 
 & 
 of the initial velocity, where $0 \leq \epsilon \ll 1$ 
 &
 %
 \\
 \hline
 (6)
 & 
 The velocity field has continuous second 
 &
 \multirow{ 2}{*}{$\dst\frac{\partial}{\partial t}
 \left(\grad \cdot \vv \right)
 = 
 \grad \cdot \left(\dst\frac{\partial \vv}{\partial t}\right)$} \\
 % 
 & 
 partial derivatives, see (\ref{eq:schwarz})
 &
 %
 \\
 \hline
\end{tabular}
\label{eq:wave_equation_assumptions}
\ee
With these additional assumptions, the {\bf elastic fluid
equations} become the {\bf wave equations}.

\begin{Hresult}
The {\bf wave equation} can be written as
\be
\boxed{
\dst\frac{\partial^2 \rho(\vx,t)}{\partial t^2} 
= 
c^2 \; \Delta \rho(\vx, t)
\;\; \Longleftrightarrow \;\;
\rho_{,tt} = c^2 \; \rho_{,ii}
}
\ee
where 
\be
 c^2 = \left.
 \dst\frac{\partial P}{\partial \rho}
 \right|_{\rho = \bar\rho_0}
\ee
and $c$ is called the {\bf sound speed} at $\bar\rho_0$.
\end{Hresult}

\begin{Hproof}
Use conservation of momentum in the current configuration
(\ref{eq:elastic_fluid_momentum})
with the 
assumptions in 
(\ref{eq:wave_equation_assumptions})
of
negligible body force, constant $\partial P / \partial \rho$,
and uniform density, 
\be
\bar\rho_0
\left[ 
\frac{\partial \vv(\vx, t)}{\partial t} 
+ 
\left[\grad \vv(\vx,t)\right]
\; \vv(\vx, t)
\right] 
= -c^2 \; \grad \rho(\vx, t).
\ee
Eliminate the $[\grad \vv] \; \vv$ term because
$\vv(\vx, t) = \epsilon \; \vv(\vx, 0) \approx \ve{0}$, and 
take the divergence,
\be
\bar\rho_0
\left[\grad \cdot \dst\frac{\partial \vv(\vx,t)}{\partial t} \right]
=
-c^2 \; \grad \cdot \grad \rho(\vx, t).
\label{eq:wave_equation_div}
\ee
Use conservation of mass from 
(\ref{eq:elastic_fluid_mass})
\be
\dst\frac{\partial \rho(\vx,t)}{\partial t} 
+
\grad \cdot
( \bar\rho_0 \; \vv(\vx,t)) 
= 
0,
\ee
and take the time derivative, 
\be
\dst\frac{\partial^2 \rho(\vx,t)}{\partial t^2} 
+
\bar\rho_0  \;
\dst\frac{\partial}{\partial t}
(\grad \cdot \vv(\vx,t))
= 
0.
\ee
Use the continuous second partial derivative assumption 
of
(\ref{eq:wave_equation_assumptions})
to rewrite the foregoing equation as
\be
\dst\frac{\partial^2 \rho(\vx, t)}{\partial t^2} 
=
-\bar\rho_0  \;
\grad \cdot
\left(\dst\frac{\partial \vv(\vx, t)}{\partial t}\right).
\ee
Finally, substitute the foregoing equation into 
(\ref{eq:wave_equation_div})
gives
\be
\dst\frac{\partial^2 \rho(\vx, t)}{\partial t^2} 
=
c^2 \; \grad \cdot \grad \rho(\vx, t)
= 
c^2 \; \Delta \rho(\vx, t),
\ee
which is the wave equation.
\end{Hproof}

\section{Incompressible, Viscid Fluid (Newtonian Fluid)}

\section{Compressible, Viscid Fluid}





